\section{凝聚 Tannaka 对偶与凝聚完备性}\label{src:7.6}
\begin{epigraph}
范畴代数的目的, 或许就是说明那些平凡的东西是平凡地平凡的.
\epigraphauthor{Peter Freyd}
\end{epigraph}
我们证明具有仿射对角态射的 Noether 代数叠的一个 Tannaka 对偶版本 (定理~\bookref{7.6.1}). 引入代数叠 $\mathcal X$ 的凝聚完备性 (定义~\bookref{7.6.5}) 后, 将证明某些有唯一闭点的商叠凝聚完备 (定理~\bookref{7.6.13}), 包括 \S\bookref{7.9.2-subsection} 对完备离散赋值环 $R$ 定义的重要例子 $\Theta_R$ 与 $\Phi_R$.

凝聚 Tannaka 对偶与凝聚完备性相结合, 使我们能将从 $\mathcal X_0$ 的各阶幂零加厚来的相容映射 $\mathcal X_n\to\mathcal Y$ 延拓为 $\mathcal X\to\mathcal Y$ (推论~\bookref{7.6.9}). 这一技巧在代数叠局部结构定理~\bookref{7.7.1} 的证明中至关重要, 也出现在许多其他论证中.

\subsection{凝聚 Tannaka 对偶}\label{src:7.6.1-subsection}
Gabriel 的经典定理 \cite{Gab62} 说, 两个 Noether 概形 $X,Y$ 同构, 当且仅当其凝聚层的抽象范畴 $\mathsf{Coh}(X)$ 与 $\mathsf{Coh}(Y)$ 等价, 即概形 $X$ 可由范畴 $\mathsf{Coh}(X)$ 恢复. 在表示论中, Saavedra Rivano 的经典 Tannaka 对偶 \cite{SR72} (另见 Deligne 与 Milne 的文章 \cite[第二章]{DMOS82}) 断言域 $\Bbbk$ 上的仿射群概形 $G$ 可由有限维表示的张量范畴 $\operatorname{Rep}^{\mathrm{fd}}(G)$ 及其遗忘函子 $\operatorname{Rep}^{\mathrm{fd}}(G)\to\mathsf{Vect}_{\Bbbk}$ 恢复.

结合这两个事实, 人们或许希望代数叠 $\mathcal X$ 可由张量范畴 $\mathsf{Coh}(\mathcal X)$ 恢复.\footnote{仅有 Abel 范畴结构不够, 例如 $\mathsf{Coh}(B(\mathbb Z/2))\cong\mathsf{Coh}(\Spec\Bbbk\amalg\Spec\Bbbk)$.} 遵循 Lurie 的精彩观察 \cite{Lur04}, 我们不但证实这一期望, 而且证明张量函子 $\mathsf{Coh}(\mathcal Y)\to\mathsf{Coh}(\mathcal X)$ 足以恢复代数叠的态射 $\mathcal X\to\mathcal Y$.
\begin{theorem}[凝聚 Tannaka 对偶]\label{src:7.6.1}
对具有仿射对角态射的 Noether 代数叠 $\mathcal X,\mathcal Y$, 函子
\begin{align}\label{src:7.6.2}
\operatorname{Mor}(\mathcal X,\mathcal Y)\longrightarrow\operatorname{Mor}^{\otimes}(\mathsf{Coh}(\mathcal Y),\mathsf{Coh}(\mathcal X)),\qquad f\longmapsto f^*
\end{align}
是范畴等价, 其中右侧表示对称幺半 Abel 范畴之间的右正合可加张量函子 $\mathsf{Coh}(\mathcal Y)\to\mathsf{Coh}(\mathcal X)$ 的范畴, 态射为张量自然变换.
\end{theorem}
\begin{remark}\label{src:7.6.3}
\newterm{对称幺半范畴}{symmetric monoidal category} 是一个范畴 $\mathcal A$, 配备双函子 $\otimes:\mathcal A\times\mathcal A\to\mathcal A$, 单位 $1\in\mathcal A$, 结合同构 $\alpha_{A,B,C}:A\otimes(B\otimes C)\xrightarrow{\sim}(A\otimes B)\otimes C$, 左右单位同构 $l_A:1\otimes A\xrightarrow{\sim}A\xrightarrow{\sim}A$ 与 $r_A:A\otimes1\to A$, 以及交换同构 $s_{A,B}:A\otimes B\cong B\otimes A$, 满足 $s_{A,B}\circ s_{B,A}=\operatorname{id}$ 及某些协调条件 \cite[\S XI.1]{Mac71}. 对称幺半 Abel 范畴之间的\newterm{张量函子}{tensor functor} $F:\mathcal A\to\mathcal B$ 是配备同构 $\Phi_{A,B}:F(A)\otimes F(B)\xrightarrow{\sim}F(A\otimes B)$ 与 $\varphi:1_{\mathcal B}\xrightarrow{\sim}F(1_{\mathcal A})$ 的函子, 这些同构与 $\alpha_{A,B,C},l_A,r_A,s_{A,B}$ 相容 \cite[\S XI.2]{Mac71}. 张量函子之间的\newterm{张量自然变换}{tensor natural transformation} 是与 $\Phi_{A,B}$ 及 $\varphi$ 相容的自然变换 \cite[\S XI.2]{Mac71}. 对称幺半 Abel 范畴 (或环 $R$ 上的对称幺半 $R$-线性 Abel 范畴) 是对称幺半 (或 $R$-线性) Abel 范畴 $\mathcal A$, 其 $\otimes$ 对每个变量可加 (或 $R$-线性). 若张量函子的底层函子可加或 $R$-线性, 也如此称该张量函子. 若 $\mathcal X,\mathcal Y$ 定义在 Noether 环 $R$ 上, 则定理~\bookref{7.6.1} 诱导等价
\begin{align*}
\operatorname{Mor}_R(\mathcal X,\mathcal Y)\xrightarrow{\sim}\operatorname{Mor}_R^{\otimes}(\mathsf{Coh}(\mathcal Y),\mathsf{Coh}(\mathcal X)),
\end{align*}
对应 $R$ 上的态射与右正合 $R$-线性张量函子.
\end{remark}
\begin{proof}
因 Noether 代数叠上每个拟凝聚层都是其凝聚子层的余极限 (命题~\bookref{7.1.15}), 每个右正合张量函子 $F:\mathsf{Coh}(\mathcal Y)\to\mathsf{Coh}(\mathcal X)$ 延拓为保持余极限的张量函子 $F:\QCoh(\mathcal Y)\to\QCoh(\mathcal X)$. 同样, 凝聚层函子之间的张量自然变换唯一延拓到拟凝聚层.

\emph{全忠实性.} 设 $f,g:\mathcal X\to\mathcal Y$. 取仿射概形的光滑展示 $p:U\to\mathcal Y$. 问题在 $\mathcal X$ 上光滑局部, 将 $\mathcal X$ 换为 $\mathcal X\times_{f,\mathcal Y,p}U$, 可假定有分解 $f:\mathcal X\xrightarrow{\widetilde f}U\xrightarrow p\mathcal Y$. 同样可假定 $g:\mathcal X\xrightarrow{\widetilde g}V\xrightarrow q\mathcal Y$, $V$ 仿射. 因 $\mathcal Y$ 的对角态射仿射, $p$ 仿射, 有等同
\begin{align*}
\operatorname{Mor}_{\mathcal Y}(\mathcal X,U)&\cong\Hom_{\mathcal O_{\mathcal Y}\text{-alg}}(p_*\mathcal O_U,f_*\mathcal O_{\mathcal X})\\
&\cong\Hom_{\mathcal O_{\mathcal X}\text{-alg}}(f^*p_*\mathcal O_U,\mathcal O_{\mathcal X}).
\end{align*}
所以 $\widetilde f,\widetilde g$ 对应于截面 $s^{\widetilde f}:f^*p_*\mathcal O_U\to\mathcal O_{\mathcal X}$ 与 $s^{\widetilde g}:g^*q_*\mathcal O_V\to\mathcal O_{\mathcal X}$. 一个 $2$-同构 $\alpha:f\to g$ 可等同于分解
\begin{align*}
\begin{tikzcd}[ampersand replacement=\&]
\mathcal X\arrow[r,"{(\widetilde f,\widetilde g,\alpha)}"]\arrow[dr,"{\operatorname{id}}"']\&U\times_{\mathcal X}V\arrow[d,"{\pi}"]\\
\&\mathcal X,
\end{tikzcd}
\end{align*}
也即 $\mathcal O_{\mathcal X}\to f^*\pi_*\mathcal O_{U\times_{\mathcal X}V}$ 的截面 $s_\alpha$. 令 $\alpha^*:f^*\to g^*$ 是 $\alpha$ 在 \bookeqref{7.6.2} 下的像, 即拉回张量自然变换, 则 $s_\alpha$ 可写为
\begin{align*}
f^*\pi_*\mathcal O_{U\times_{\mathcal X}V}&\cong f^*(p_*\mathcal O_U)\otimes f^*(q_*\mathcal O_V)\\
&\xrightarrow{\operatorname{id}\otimes\alpha^*_{q_*\mathcal O_V}}f^*(p_*\mathcal O_U)\otimes g^*(q_*\mathcal O_V)\xrightarrow{s^{\widetilde f}\otimes s^{\widetilde g}}\mathcal O_S.
\end{align*}
若 $2$-同构 $\alpha,\alpha':f\to g$ 满足 $\alpha^*=\alpha'^*$, 则 $\alpha^*_{q_*\mathcal O_V}=\alpha'^*_{q_*\mathcal O_V}$, 所以 $s_\alpha=s_{\alpha'}$, $\alpha=\alpha'$, 这给出忠实性. 对于完全性, 设 $\beta:f^*\to g^*$ 是张量自然变换. 则 $\operatorname{id}\otimes\beta_{q_*\mathcal O_V}$ 定义截面 $f^*\pi_*\mathcal O_{U\times_{\mathcal X}V}\to\mathcal O_S$, 从而给出 $2$-同构 $\alpha:f\to g$, 满足 $\beta_{q_*\mathcal O_V}=\alpha^*_{q_*\mathcal O_V}$. 为证明每个 $E\in\QCoh(\mathcal Y)$ 满足 $\beta_E=\alpha_E^*$, 注意分解 $g=q\circ\widetilde g$ 给出 $g^*E\to g^*(q_*q^*E)$ 的分裂. 因 $f^*,g^*$ 与直和交换, 只需假定 $E=q_*G$, $G\in\QCoh(V)$. 将 $G=\operatorname{colim}(\mathcal O_V^{\oplus I}\to\mathcal O_V^{\oplus J})$ 写为自由 $\mathcal O_V$-模的余极限. 因 $f^*,g^*$ 与余极限交换且 $q_*$ 正合, 得到 $\beta_{q_*G}=\alpha^*_{q_*G}$.

\emph{本质满性 (仿射情形).} 设 $F:\QCoh(\mathcal Y)\to\QCoh(\mathcal X)$ 是保持余极限的张量函子. 若 $\mathcal X=\Spec A$, $\mathcal Y=\Spec B$ 是 Noether 仿射概形, 有映射
\begin{align*}
\phi:B\cong\End(\mathcal O_{\mathcal Y})\xrightarrow F\End(\mathcal O_{\mathcal X})\cong A.
\end{align*}
断言 $\phi$ 是环同态, 且对 $N\in\Mod_B$ 有函子性同构 $F(N)=N\otimes_BA$. 对 $b,b'\in B$, 考虑交换图
\begin{align*}
\begin{tikzcd}[ampersand replacement=\&,column sep=small]
\mathcal O_{\mathcal Y}\otimes\mathcal O_{\mathcal Y}\arrow[r]\arrow[d,"{b\otimes b'}"']\&\mathcal O_{\mathcal Y}\arrow[d,"{bb'}"]\\
\mathcal O_{\mathcal Y}\otimes\mathcal O_{\mathcal Y}\arrow[r]\&\mathcal O_{\mathcal Y}
\end{tikzcd}
\quad\xrightarrow{F}\quad
\begin{tikzcd}[ampersand replacement=\&,column sep=small]
A\otimes A\arrow[r]\arrow[d,"{\phi(b)\otimes\phi(b')}"']\&A\arrow[d,"{\phi(bb')}"]\\
A\otimes A\arrow[r]\&A,
\end{tikzcd}
\end{align*}
其中横箭头为乘法. $F$ 保持张量积使右图交换, 故 $\phi(b)\phi(b')=\phi(bb')$. 对 $B$-模 $N$, 取自由展示 $B^{\oplus J}\to B^{\oplus I}\to N\to0$. 因 $F$ 与 $-\otimes_BA$ 都右正合且保持直和, 应用于此展示后, 两者都是 $A^{\oplus J}\to A^{\oplus I}$ 的余核, 得到 $F(N)\cong N\otimes_BA$. 同样可检验其函子性.

\emph{归约到 $\mathcal X$ 仿射.} 取仿射概形的光滑展示 $g:U\to\mathcal X$, 考虑图
\begin{align*}
\begin{tikzcd}[ampersand replacement=\&]
U\times_{\mathcal X}U\arrow[d,shift left,"{p_1}"]\arrow[d,shift right,"{p_2}"']\&\\
U\arrow[d,"{g}"']\arrow[dr,dotted,"{g^*\circ F}"]\&\\
\mathcal X\arrow[r,dotted,"{F}"']\&\mathcal Y,
\end{tikzcd}
\end{align*}
其中虚箭头 $\mathcal X\dashrightarrow\mathcal Y$ 表示反方向的张量函子 $F:\QCoh(\mathcal Y)\to\QCoh(\mathcal X)$. 若 $U$ 仿射时结论成立, 则有态射 $h:U\to\mathcal Y$ 与函子的同构 $h\xrightarrow{\sim}g^*\circ F$. 由全忠实性, 有满足余循环条件的同构 $p_1\circ h\xrightarrow{\sim}p_2\circ h$, 故光滑下降给出唯一态射 $f:\mathcal X\to\mathcal Y$, 使 $F\simeq f^*$.

\emph{归约到 $\mathcal Y$ 仿射.} 令 $\mathcal X=\Spec A$, 取光滑展示 $q:V=\Spec C\to\mathcal Y$. 因 $\mathcal Y$ 的对角态射仿射, $q$ 仿射. 定义 $B:=F(q_*\mathcal O_V)$, 它是 $A$-代数, 因为 $q_*\mathcal O_V$ 是 $\mathcal O_V$-代数. 考虑图
\begin{align*}
\begin{tikzcd}[ampersand replacement=\&]
\Spec B\arrow[r,dotted,"{F'}"]\arrow[d]\&V=\Spec C\arrow[d,"{q}"]\\
\Spec A=\mathcal X\arrow[r,dotted,"{F}"']\&\mathcal Y,
\end{tikzcd}
\end{align*}
其中 $F':\Mod_C\to\Mod_B$ 是将 $M$ 送到 $F(q_*\widetilde M)$ 的右正合张量函子. 后者是 $B=F(q_*\mathcal O_V)$-模, 因 $q_*\widetilde M$ 是 $q_*\mathcal O_V$-模. 仿射情形说明 $F'$ 由态射 $f':\Spec B\to\Spec C$ 诱导. 图可扩充为
\begin{align*}
\begin{tikzcd}[ampersand replacement=\&]
\Spec(B\otimes_AB)\arrow[r,"{f''}"]\arrow[d,shift left]\arrow[d,shift right]\&V\times_{\mathcal Y}V\arrow[d,shift left]\arrow[d,shift right]\\
\Spec B\arrow[r,"{f'}"]\arrow[d]\&V=\Spec C\arrow[d,"{q}"]\\
\Spec A=\mathcal X\arrow[r,dotted,"{F}"']\&\mathcal Y.
\end{tikzcd}
\end{align*}
因 $q$ 仿射, $V\times_{\mathcal Y}V$ 仿射, 上方的方块对任一对投影均笛卡尔.

若能证明 $A\to B$ 忠实平坦, 就可由仿射情形的全忠实性使 $f'$ 下降为所需 $f:\mathcal X\to\mathcal Y$. 直接检验似乎困难, 但已知 $B\rightrightarrows B\otimes_AB$ 忠实平坦, 因为它们对应于光滑映射 $V\times_{\mathcal Y}V\rightrightarrows V$ 的基变换. 我们转而证明 $A\to B$ 泛单射. 忠实平坦性沿泛单射映射下降 (命题~\bookref{A.2.24}(4)), 故此即蕴含忠实平坦性.

\emph{$A\to B$ 的泛单射性.} 回忆定义~\bookref{A.2.22}, $A$-模单射在与每个 $A$-模张量后仍单射时称为泛单射. 由命题~\bookref{A.2.24}(3), 它在忠实平坦态射下是局部性质, 从而可推广到代数叠上拟凝聚层的态射. 因 $q$ 忠实平坦, $\mathcal O_{\mathcal Y}\to q_*\mathcal O_V$ 泛单射 (命题~\bookref{A.2.24}(1)). 写 $q_*\mathcal O_V=\operatorname{colim}Q_i$ 为凝聚子层的余极限 (命题~\bookref{7.1.15}), 并可假定每个 $Q_i$ 包含 $\mathcal O_{\mathcal Y}$ 的像. 则 $\mathcal O_{\mathcal Y}\to Q_i$ 也泛单射. 因 $Q_i$ 凝聚, 该映射在 $\mathcal Y$ 上光滑局部是分裂单射 (命题~\bookref{A.2.24}(2)). 应用 $F$ 并利用其保持余极限, 得到分解
\begin{align*}
\begin{tikzcd}[ampersand replacement=\&,column sep=large]
\&F(Q_i)\arrow[d]\\
A=F(\mathcal O_{\mathcal Y})\arrow[ur]\arrow[r]\&B=F(q_*\mathcal O_V)=\operatorname{colim}F(Q_i).
\end{tikzcd}
\end{align*}
只需证明 $A\to F(Q_i)$ 泛单射. 实际上它分裂. 因 $\mathcal O_{\mathcal Y}\to Q_i$ 光滑局部分裂, 对偶映射 $Q_i^\vee\to\mathcal O_{\mathcal Y}^\vee=\mathcal O_{\mathcal Y}$ 满射. 由 $F$ 右正合, 得到满射 $F(Q_i^\vee)\to F(\mathcal O_{\mathcal Y})=A$, 可取映到 $1$ 的元素 $\lambda$. 自然映射 $F(Q_i^\vee)\to F(Q_i)^\vee$ 将 $\lambda$ 送到 $F(Q_i)\to A$, 它给出 $A\to F(Q_i)$ 的分裂. 另见 \cite{Lur04}, \cite{HR19b}, \cite{BHL17} 与 \cite[Tag 0GRR]{SP}.
\end{proof}
\begin{remark}[与经典 Tannaka 对偶的关系]\label{src:7.6.4}
若 $G$ 是域 $\Bbbk$ 上的仿射群概形, 则有限维表示的范畴 $\mathcal C=\operatorname{Rep}^{\mathrm{fd}}(G)$ 是对称幺半 $\Bbbk$-线性范畴, 并有张量函子 $\omega:\operatorname{Rep}^{\mathrm{fd}}(G)\to\mathsf{Vect}_{\Bbbk}$. 对 $\Bbbk$-代数 $R$, 定义复合
\begin{align*}
\omega_R:\operatorname{Rep}^{\mathrm{fd}}(G)\to\mathsf{Vect}_{\Bbbk}\xrightarrow{-\otimes_{\Bbbk}R}\Mod_R,
\end{align*}
令 $\operatorname{Aut}^{\otimes}(\omega_R)$ 为 $\omega_R$ 的张量自然同构群. 则 $G$ 可恢复为仿射 $\Bbbk$-概形上的函子 $\underline{\operatorname{Aut}}^{\otimes}(\omega)$, 它将 $R$ 送到 $\operatorname{Aut}^{\otimes}(\omega_R)$ \cite[II.2.8]{DMOS82}. 另一方面, 凝聚 Tannaka 对偶~\bookref{7.6.1} 断言每个 Noether $\Bbbk$-代数 $R$ 有范畴等价
\begin{align*}
\operatorname{Mor}_{\Bbbk}(\Spec R,BG)\xrightarrow{\sim}\operatorname{Mor}^{\otimes}(\operatorname{Rep}(G)^{\mathrm{fd}},\Mod_R).
\end{align*}
由此 $\operatorname{Rep}(G)^{\mathrm{fd}}$ 决定 $BG$. 为恢复 $G$, 纤维函子 $\omega$ 对应于态射 $p:\Spec\Bbbk\to BG$, 且 $G=\underline{\operatorname{Aut}}_{\Bbbk}(p)$. 例如, 考虑同维非退化二次型 $q,q'$ 的正交群 $O(q),O(q')$. 虽然这些群未必同构, 却有 $\operatorname{Rep}(O(q))\cong\operatorname{Rep}(O(q'))$. 此时两个映射 $\Spec\Bbbk\to BO(q)$ 与 $\Spec\Bbbk\to BO(q')$ 在同一范畴上定义两个不同纤维函子.

经典版本还给出数据 $(\mathcal C,\omega)$ 同构于群概形表示范畴的条件. 若对每个 $X\in\mathcal C$, 存在对偶 $X^\vee\in\mathcal C$, 即 $X^\vee\otimes-:\mathcal C\to\mathcal C$ 为 $X\otimes-$ 的右伴随, 则称 $\mathcal C$ \newterm{刚性的}{rigid}. 若 $\mathcal C$ 是刚性对称幺半 $\Bbbk$-线性 Abel 范畴, $\End(1)=\Bbbk$, 且 $\omega:\mathcal C\to\mathsf{Vect}_{\Bbbk}$ 是正合忠实的 $\Bbbk$-线性张量函子, 则 $\operatorname{Aut}^{\otimes}(\omega)$ 可由仿射 $\Bbbk$-群概形 $G$ 表示, 并有张量等价 $\mathcal C\cong\operatorname{Rep}^{\mathrm{fd}}(G)$, 在此等价下 $\omega$ 对应于遗忘函子 \cite[II.2.11]{DMOS82}. 此外 $G$ 在 $\Bbbk$ 上有限型当且仅当 $\mathcal C$ 有张量生成元.
\end{remark}

\subsection{凝聚完备性}\label{src:7.6.2-subsection}
凝聚 Tannaka 对偶与凝聚完备性的结合 (推论~\bookref{7.6.9}) 将成为我们手中一个有价值的新工具.
\begin{definition}\label{src:7.6.5}
称 Noether 代数叠 $\mathcal X$ \newterm{沿闭子叠 $\mathcal X_0$ 凝聚完备}{coherently complete along a closed substack}, 若自然函子
\begin{align*}
\mathsf{Coh}(\mathcal X)\longrightarrow\varprojlim\mathsf{Coh}(\mathcal X_n),\qquad F\longmapsto(F_n)
\end{align*}
是范畴等价, 其中 $\mathcal X_n$ 是 $\mathcal X_0$ 的第 $n$ 阶幂零加厚, $F_n$ 为 $F$ 到 $\mathcal X_n$ 的拉回.
\end{definition}
\begin{remark}\label{src:7.6.6}
若凝聚理想层 $\mathcal I\subseteq\mathcal O_{\mathcal X}$ 定义 $\mathcal X_0$, 则 $\mathcal X_n$ 由 $\mathcal I^{n+1}$ 定义. 记自然包含为 $i_n:\mathcal X_n\hookrightarrow\mathcal X_{n+1}$, 则 $\varprojlim\mathsf{Coh}(\mathcal X_n)$ 中的对象对应于凝聚层序列 $F_n\in\mathsf{Coh}(\mathcal X_n)$, 以及诱导同构 $F_n\to i_n^*F_{n+1}$ 的映射 $\alpha_n:i_{n,*}F_n\to F_{n+1}$. 态射 $(F_n,\alpha_n)\to(F'_n,\alpha'_n)$ 是满足 $\phi_{n+1}\circ\alpha_n=\alpha'_{n+1}\circ i_{n,*}\phi_n$ 的映射序列 $\phi_n:F_n\to F'_n$.
\end{remark}
\begin{example}\label{src:7.6.7}
若 $(R,\mathfrak m)$ 是完备 Noether 局部环, 则 Artin--Rees 引理~\bookref{B.5.4} 说明 $\Spec R$ 沿 $\Spec R/\mathfrak m$ 凝聚完备. 对 Noether $I$-进完备环 $R=\varprojlim R/I^n$, 同样成立.
\end{example}
\begin{example}\label{src:7.6.8}
Grothendieck 存在定理~\bookref{C.5.3} 断言, 若 $X$ 是完备局部环 $(R,\mathfrak m)$ 上的固有概形, $X_0=X\times_RR/\mathfrak m$, 则 $X$ 沿 $X_0$ 凝聚完备. 若 $\mathcal X$ 是 $\Spec R$ 上的固有 Deligne--Mumford 叠, 同样成立 (习题~\bookref{5.6.13}(b)). 对任意 $I$-进完备 Noether 环上的固有代数叠也成立. 见 \cite[定理 1.4]{Ols05} 或 \cite[定理 4.1]{Con05a}.
\end{example}
\begin{corollary}[凝聚 Tannaka 对偶]\label{src:7.6.9}
设 $\mathcal X,\mathcal Y$ 是具有仿射对角态射的 Noether 代数叠, $\mathcal X$ 沿 $\mathcal X_0$ 凝聚完备. 则有范畴等价
\begin{align*}
\operatorname{Mor}(\mathcal X,\mathcal Y)\longrightarrow\varprojlim\operatorname{Mor}(\mathcal X_n,\mathcal Y),\qquad f\longmapsto(f_n),
\end{align*}
其中 $f_n:\mathcal X_n\to Y$ 是 $f$ 到 $\mathcal X_0$ 的第 $n$ 阶幂零加厚 $\mathcal X_n$ 的限制.
\end{corollary}
\begin{proof}
由以下等价得到:
\begin{align*}
\operatorname{Mor}(\mathcal X,\mathcal Y)&\simeq\operatorname{Mor}^{\otimes}(\mathsf{Coh}(\mathcal Y),\mathsf{Coh}(\mathcal X))&&\text{(凝聚 Tannaka 对偶)}\\
&\simeq\operatorname{Mor}^{\otimes}(\mathsf{Coh}(\mathcal Y),\varprojlim\mathsf{Coh}(\mathcal X_n))&&\text{(凝聚完备性)}\\
&\simeq\varprojlim\operatorname{Mor}^{\otimes}(\mathsf{Coh}(\mathcal Y),\mathsf{Coh}(\mathcal X_n))\\
&\simeq\varprojlim\operatorname{Mor}(\mathcal X_n,\mathcal Y)&&\text{(凝聚 Tannaka 对偶)}.
\end{align*}
\end{proof}
\begin{remark}\label{src:7.6.10}
若 $\mathcal X,\mathcal Y$ 定义在 Noether 环 $R$ 上, 则有等价 $\operatorname{Mor}_R(\mathcal X,\mathcal Y)\to\varprojlim\operatorname{Mor}_R(\mathcal X_n,\mathcal Y)$. 这由 $R$ 上态射的范畴与右正合 $R$-线性张量函子范畴之间的 Tannaka 对偶 (注记~\bookref{7.6.3}) 同样得到.
\end{remark}
\begin{exercise}\label{src:7.6.11}
当 $(R,\mathfrak m)$ 是完备 Noether 局部环时, 直接证明 $\operatorname{Mor}(\Spec R,\mathcal Y)\to\varprojlim\operatorname{Mor}(\Spec R/\mathfrak m^n,\mathcal Y)$ 是等价.

提示: 用命题~\bookref{5.3.4} 将 $\Spec R/\mathfrak m\to\mathcal Y$ 提升到光滑展示, 再用无穷小提升判据~\bookref{4.7.1} 归约到 $\mathcal Y$ 为概形的情形.
\end{exercise}
\begin{exercise}\label{src:7.6.12}
设仿射代数群 $G$ 作用于分离 Noether $\Bbbk$-代数空间 $W$, $W_0\subseteq W$ 是 $G$-不变闭子空间, $W_n$ 为其第 $n$ 阶幂零加厚. 假定 $[W/G]$ 沿 $[W_0/G]$ 凝聚完备. 对每个配备 $G$-作用且具有仿射对角态射的 Noether $\Bbbk$-代数空间 $X$, 等变映射的自然映射
\begin{align*}
\operatorname{Mor}^G(W,X)\longrightarrow\varprojlim_n\operatorname{Mor}^G(W_n,X)
\end{align*}
双射.

提示: $G$-等变映射 $W\to X$ 对应于 $BG$ 上的态射 $[W/G]\to[X/G]$, 用此归约到推论~\bookref{7.6.9}.
\end{exercise}

\subsection{商叠的凝聚完备性}\label{src:7.6.3-subsection}
我们将应用的凝聚完备性结果 (例如用于代数叠局部结构定理~\bookref{7.7.1} 的证明) 如下:
\begin{theorem}\label{src:7.6.13}
设 $\Bbbk$ 是代数闭域, $R$ 是剩余域为 $\Bbbk$ 的完备 Noether 局部 $\Bbbk$-代数. 设 $\Bbbk$ 上的线性约化群 $G$ 作用于 $R$ 上有限型的仿射概形 $\Spec A$. 假定 $A^G=R$, 且存在 $G$-不动 $\Bbbk$-点 $x\in\Spec A$. 则 $[\Spec A/G]$ 沿 $x$ 定义的闭子叠 $BG$ 凝聚完备.
\end{theorem}
\begin{example}\label{src:7.6.14}
若 $\mathbb G_m$ 对角作用于 $\mathbb A^r$, 则 $[\mathbb A^r/\mathbb G_m]$ 沿原点 $B\mathbb G_m$ 凝聚完备. 即 $\Bbbk[x_1,\ldots,x_r]$ 上的 $\mathbb G_m$-等变模等价于 $\Bbbk[x_1,\ldots,x_r]/(x_1,\ldots,x_r)^{n+1}$ 上 $\mathbb G_m$-等变模的相容族.
\end{example}
\begin{remark}\label{src:7.6.15}
考虑图
\begin{align*}
\begin{tikzcd}[ampersand replacement=\&,column sep=large]
BG\arrow[r,hook]\&{[\Spec A/G]\times_{A^G}\Bbbk}\arrow[r,hook]\arrow[d]\&{[\Spec A/G]}\arrow[d]\\
\&\Spec\Bbbk\arrow[r,hook]\&\Spec A^G.
\end{tikzcd}
\end{align*}
上述定理形式上蕴含 $[\Spec A/G]$ 对纤维 $[\Spec A/G]\times_{A^G}\Bbbk$ 也凝聚完备. 这一版本类似 Grothendieck 存在定理 (例~\bookref{7.6.8}), 但沿 $BG$ 的凝聚完备性强得多. 例如 $[\mathbb A^n/\mathbb G_m]\to\Spec\Bbbk$ 的纤维就是整个叠.
\end{remark}
\begin{proof}[定理~\bookref{7.6.13} 的证明]
需证明 $\mathsf{Coh}(\mathcal X)\to\varprojlim\mathsf{Coh}(\mathcal X_n)$ 是范畴等价, 其中 $\mathcal X=[\Spec A/G]$, $\mathcal X_n$ 是 $x$ 处剩余 gerbe 的包含 $BG\hookrightarrow\mathcal X$ 的第 $n$ 阶幂零加厚.

\emph{全忠实性.} 设 $F,F'$ 是凝聚 $\mathcal O_{\mathcal X}$-模, 其到 $\mathcal X_n$ 的限制分别为 $F_n,F'_n$. 要证明 $\Hom_{\mathcal O_{\mathcal X}}(F,F')\to\varprojlim\Hom_{\mathcal O_{\mathcal X}}(F_n,F'_n)$ 双射. 因 $\mathcal X$ 具有消解性质 (命题~\bookref{7.2.10}), 可取向量丛消解 $F_2\to F_1\to F\to0$. 这诱导正合行的图
\begin{align*}
\begin{tikzcd}[ampersand replacement=\&,column sep=small]
0\arrow[r]\&\Hom(F,F')\arrow[r]\arrow[d]\&\Hom(F_1,F')\arrow[r]\arrow[d]\&\Hom(F_2,F')\arrow[d]\\
0\arrow[r]\&\varprojlim\Hom(F_n,F'_n)\arrow[r]\&\varprojlim\Hom(F_{1,n},F'_n)\arrow[r]\&\varprojlim\Hom(F_{2,n},F'_n),
\end{tikzcd}
\end{align*}
其中 $\Hom$ 均为 $\mathcal O_{\mathcal X}$-模的态射. 因此只需在 $F$ 为向量丛时证明全忠实性. 此时
\begin{align*}
\Hom_{\mathcal O_{\mathcal X}}(F,F')&=\Hom_{\mathcal O_{\mathcal X}}(\mathcal O_{\mathcal X},F^\vee\otimes F'),\\
\Hom_{\mathcal O_{\mathcal X}}(F_n,F'_n)&=\Hom_{\mathcal O_{\mathcal X}}(\mathcal O_{\mathcal X_n},F_n^\vee\otimes F'_n).
\end{align*}
所以还可假定 $F=\mathcal O_{\mathcal X}$, 归约为证明
\begin{align}\label{src:7.6.16}
\Gamma(\mathcal X,F')\longrightarrow\varprojlim\Gamma(\mathcal X_n,F'_n)
\end{align}
是同构. 写 $F'=\widetilde M$, $M$ 是带有 $G$-作用的有限生成 $A$-模, 令 $\mathfrak m\subseteq A$ 是 $x$ 的极大理想. 因 $G$ 线性约化, $\Gamma(\mathcal X_n,F'_n)=M^G/(\mathfrak m^nM)^G$. 因此需验证
\begin{align}\label{src:7.6.17}
M^G\longrightarrow\varprojlim M^G/(\mathfrak m^nM)^G
\end{align}
是同构. 先证明
\begin{align}\label{src:7.6.18}
\bigcap_{n\ge0}(\mathfrak m^nM)^G=0,
\end{align}
即 \bookeqref{7.6.17} 单射. 令 $N:=\bigcap_{n\ge0}\mathfrak m^nM$. 对 $N\subseteq M$ 应用 Artin--Rees 引理~\bookref{B.5.4}, 有整数 $c$ 使所有 $n\ge c$ 满足 $\mathfrak m^nM\cap N=\mathfrak m^{n-c}(\mathfrak m^cM\cap N)$. 取 $n=c+1$, 得到 $N=\mathfrak mN$, 故 $N\otimes_AA/\mathfrak m=0$. $N$ 的支撑是 $\Spec A$ 中不包含 $x$ 的闭 $G$-不变子概形, 因而 $N=0$.

因 $G$ 线性约化, $M^G$ 是有限生成 $A^G$-模 (推论~\bookref{7.5.13}(3)). 接着在 $A^G$ 为 Artin 环时证明 \bookeqref{7.6.17}. 此时 $((\mathfrak m^nM)^G)$ 是 Artin $A^G$-模序列, 自动满足 Mittag--Leffler 条件. 对 $0\to(\mathfrak m^nM)^G\to M^G\to M^G/(\mathfrak m^nM)^G\to0$ 取逆极限并用 \bookeqref{7.6.18}, 得到正合列
\begin{align*}
0\to0\to M^G\to\varprojlim M^G/(\mathfrak m^nM)^G\to0,
\end{align*}
故 \bookeqref{7.6.17} 为同构. 一般情形, 令 $J=\mathfrak m^GA\subseteq A$, 注意
\begin{align}\label{src:7.6.19}
M^G\cong\varprojlim M^G/(\mathfrak m^G)^nM^G\cong\varprojlim(M/J^nM)^G,
\end{align}
因为 $G$ 线性约化. 对每个 $n$, Artin 情形给出
\begin{align}\label{src:7.6.20}
(M/J^nM)^G\cong\varprojlim_l M^G/((J^n+\mathfrak m^l)M)^G.
\end{align}
结合 \bookeqref{7.6.19}, \bookeqref{7.6.20} 及 $n\ge l$ 时 $J^n\subseteq\mathfrak m^l$, 得到
\begin{align*}
M^G&\cong\varprojlim_n(M/J^nM)^G\\
&\cong\varprojlim_n\varprojlim_l M^G/((J^n+\mathfrak m^l)M)^G\\
&\cong\varprojlim_l M^G/(\mathfrak m^lM)^G.
\end{align*}

\emph{本质满性 (显式构造).} $G$ 线性约化, 所以 $[\Spec A/G]$ 上每个凝聚层 $F=\widetilde M$ 分解为直和
\begin{align}\label{src:7.6.21}
M=\bigoplus_{\rho\in\Gamma}M^{(\rho)},
\end{align}
其中 $\Gamma$ 为 $G$ 的不可约表示同构类的集合, $M^{(\rho)}$ 是对应于 $\rho$ 的\newterm{等型分量}{isotypic component}. 具体地, 若 $W_\rho$ 表示对应于 $\rho$ 的不可约表示, 则 $M^{(\rho)}=\Hom_{\Bbbk}^G(W_\rho,M)\otimes W_\rho$. 分解 \bookeqref{7.6.21} 与 $M$ 的 $A$-模结构及分解 $A=\bigoplus_{\rho\in\Gamma}A^{(\rho)}$ 相容.

若 $F=\widetilde M\in\mathsf{Coh}(\mathcal X)$ 的限制为 $M_n=M/\mathfrak m^{n+1}M$, 对 $M\otimes W_\rho^\vee$ 应用 \bookeqref{7.6.17}, 得到 $M^{(\rho)}=\varprojlim M_n^{(\rho)}$. 定理~\bookref{7.5.11}(3) 还说明 $M^{(\rho)}$ 是有限生成 $A^G$-模. 特别地, $A^{(\rho)}=\varprojlim(A/\mathfrak m^{n+1})^{(\rho)}$ 是有限生成 $A^G$-模.

于是, 若 $F_n=\widetilde M_n$ 是凝聚 $\mathcal O_{\mathcal X_n}$-模的相容系统, $M_n=\bigoplus_\rho M_n^{(\rho)}$, 可定义
\begin{align}\label{src:7.6.22}
M^{(\rho)}:=\varprojlim M_n^{(\rho)},\qquad M:=\bigoplus_{\rho\in\Gamma}M^{(\rho)}.
\end{align}
为证明 $M$ 是带 $G$-作用的 $A$-模, 取不可约表示 $\rho,\gamma\in\Gamma$, 令 $\Lambda\subseteq\Gamma$ 为 $W_\rho\otimes W_\gamma$ 中出现的非零不可约表示的有限集合. 对映射 $A_n^{(\rho)}\otimes_{(A/\mathfrak m^{n+1})^G}M_n^{(\gamma)}\to\bigoplus_{\lambda\in\Lambda}M_n^{(\lambda)}$ 取极限, 定义乘法
\begin{align*}
A^{(\rho)}\otimes_AM^{(\gamma)}&\longrightarrow\varprojlim\bigl(A_n^{(\rho)}\otimes_{(A/\mathfrak m^{n+1})^G}M_n^{(\gamma)}\bigr)\\
&\longrightarrow\varprojlim\Bigl(\bigoplus_{\lambda\in\Lambda}M_n^{(\lambda)}\Bigr)\cong\bigoplus_{\lambda\in\Lambda}M^{(\lambda)}.
\end{align*}
由构造, $M/\mathfrak m^{n+1}M\cong M_n$.

只需证明 \bookeqref{7.6.22} 的 $A$-模 $M$ 有限生成. $\mathcal X_0=BG$ 上的凝聚层 $F_0=\widetilde M_0$ 是有限维 $G$-表示, 可考虑凝聚 $\mathcal O_{\mathcal X}$-模 $F_0\otimes_{\mathcal O_{\mathcal X_0}}\mathcal O_{\mathcal X}$, 等价地为带自然 $G$-作用的 $A$-模 $M_0\otimes_{\Bbbk}A$. 因 $\mathcal X$ 上同调仿射, 函子
\begin{align*}
\Hom_{\mathcal O_{\mathcal X}}(F_0\otimes_{\mathcal O_{\mathcal X_0}}\mathcal O_{\mathcal X},-)
=\Gamma\bigl(\mathcal X,(F_0^\vee\otimes_{\mathcal O_{\mathcal X_0}}\mathcal O_{\mathcal X})\otimes_{\mathcal O_{\mathcal X}}-\bigr)
\end{align*}
正合. 应用于满射 $M\to M_0$, 得到映射
\begin{align}\label{src:7.6.23}
M_0\otimes_{\Bbbk}A\longrightarrow M,
\end{align}
要证明它满射. 已知限制 $M_0\otimes_{\Bbbk}(A/\mathfrak m^{n+1})\to M_n$ 满射, 因其余核是 $\mathcal X_n$ 上不支撑在唯一闭点处的凝聚模.

如前, 先设 $A^G$ 为 Artin 环. 因 $A^{(\rho)}\xrightarrow{\sim}\varprojlim(A/\mathfrak m^n)^{(\rho)}$ 是有限生成 $A^G$-模, 当 $n\gg0$ 时 $(A/\mathfrak m^n)^{(\rho)}$ 稳定为 $A^{(\rho)}$. 因 \bookeqref{7.6.23} 诱导满射 $M_0\otimes_{\Bbbk}(A/\mathfrak m^{n+1})\to M_n$, 当 $n\gg0$ 时 $M_n^{(\rho)}$ 稳定为 $M_\infty^{(\rho)}$, $M=\bigoplus_\rho M_\infty^{(\rho)}$ 有限生成. 一般地, 令 $X_m=\Spec A^G/(\mathfrak m\cap A^G)^{m+1}$, 考虑笛卡尔图
\begin{align*}
\begin{tikzcd}[ampersand replacement=\&,column sep=small]
\mathcal X\times_XX_m\arrow[r,hook,"{i_m}"]\arrow[d,"{\pi_m}"']\&\mathcal X\times_XX_{m+1}\arrow[r,hook]\arrow[d,"{\pi_{m+1}}"]\&\cdots\arrow[r,hook]\&\mathcal X\arrow[d,"{\pi}"]\\
X_m\arrow[r,hook,"{j_m}"']\&X_{m+1}\arrow[r,hook]\&\cdots\arrow[r,hook]\&X.
\end{tikzcd}
\end{align*}
对每个 $m$, 考虑 $\mathcal X_0\hookrightarrow\mathcal X\times_XX_m$ 的第 $n$ 阶幂零加厚 $\mathcal Z_{m,n}$, 它是 $\mathcal X_n$ 的闭子叠. 因 $X_m$ 是 Artin 环的谱, 限制 $F_n|_{\mathcal Z_{m,n}}$ 延拓为 $\mathcal X\times_XX_m$ 上的凝聚层 $H_m=\widetilde N_m$. 此外 $H_m$ 与 $H_{m+1}$ 到该叠的限制有典范同构. 引理~\bookref{7.5.24}(4) 说明拟凝聚层上的伴随态射 $j_m^*\pi_{m+1,*}\xrightarrow{\sim}\pi_{m,*}i_m^*$ 是同构. 因而
$N_{m+1}^{(\rho)}=\Gamma(\mathcal X\times_XX_{m+1},H_{m+1}\otimes W_\rho^\vee)$ 限制为 $N_m^{(\rho)}$, 且 $M^{(\rho)}=\varprojlim N_m^{(\rho)}$ 是有限生成 $A^G$-模. 由 Nakayama 引理, \bookeqref{7.6.23} 在每个 $\rho$-等型分量上满射, 故满射.

\emph{本质满性 (间接构造).} 另一论证是利用 $\mathcal X$ 的消解性质 (命题~\bookref{7.2.10}) 取向量丛的满射 $E\to F_0$. 因每个 $F_{n+1}\to F_n$ 满射, 且 $\Hom_{\mathcal O_{\mathcal X}}(E,-)=\Gamma(\mathcal X,E^\vee\otimes_{\mathcal O_{\mathcal X}}-)$ 正合, 可将 $E\to F_0$ 提升为相容映射 $E\to F_n$, 它们都满射 (Nakayama 引理). 序列 $(\ker(E_n\to F_n))$ 未必是凝聚层的进系统, 因为 $\ker(E_{n+1}\to F_{n+1})$ 到 $\mathcal X_n$ 的限制未必是 $\ker(E_n\to F_n)$. 但可如下修正: 对每个 $l\ge m\ge n$, 当 $l\gg m$ 时 $\ker(E_l\to F_l)$ 在 $E_m$ 中的像稳定为 $K'_m$, 当 $m\gg n$ 时 $K'_m/\mathfrak m^{n+1}K'_m$ 稳定为 $K_n$ (另见 \cite[Tag 087X]{SP}). 则 $(K_n)\in\varprojlim\mathsf{Coh}(\mathcal X_n)$ 是进序列. 重复构造, 可找到 $\mathcal X$ 上的向量丛 $E'$ 及相容满射 $E'\to K_n$. 全忠实性给出延拓 $E'_n\to E_n$ 的态射 $E'\to E$. 则 $\operatorname{coker}(E'\to E)$ 是延拓 $(F_n)$ 的凝聚 $\mathcal O_X$-模. 另见 \cite[定理 1.3]{AHR20} 与 \cite[定理 1.6]{AHR19}.
\end{proof}
\begin{exercise}\label{src:7.6.24}
若 $S$ 是 Noether 仿射概形, 证明 $[\mathbb A^1/\mathbb G_m]_S$ 沿 $B\mathbb G_{m,S}$ 凝聚完备.
\end{exercise}
